\section{An Exchanger}

We now use |Condition|s to implement an exchanger, following the trait in
Figure~\ref{fig:ExchangerT}.  The first thread to arrive deposits its value in
an object variable, and waits for the second thread.  The second thread, swaps
the deposited value with its own, and then signals to the first thread.
However, we need to ensure that there is no interference from other threads.
If we are not careful, a third thread could take the second thread's value
before the first thread takes it: this would not be a correct exchange.

In order to coordinate the exchange, we will use a variable |stage| that
records which stage has been reached.  This variable will store one of the
following values.
%
\begin{description}
\item[{\scalastyle Empty}:] No thread is currently exchanging;

\item[{\scalashape Filled}:] The first thread has deposited its value;

\item[{\scalashape Exchanging}:] An exchange is underway: the second thread
  has taken the first thread's value, and swapped it with its own value.
\end{description}
%
In particular, new threads will wait when $\sm{stage} = \sm{Exchanging}$, to
prevent them interfering with the current exchange.

An implementation is in Figure~\ref{fig:MonitorExchanger}.  (The trait
|Exchanger1T| extends |ExchangerT| by giving a default no-op implementation to
the |shutdown| operation.)  This uses a variable |slot|, which hods the first
thread's value when $\sm{stage} = \sm{Filled}$, holds the second thread's
value when $\sm{stage} = \sm{Exchanging}$, and is invalid when $\sm{stage} =
\sm{Empty}$. 


%%%%%

\begin{figure}
\begin{scala}
class MonitorExchanger[A] extends tacp.clientServer.Exchanger1T[A]{  
  private val Empty = 0
  private val Filled = 1
  private val Exchanging = 2
  private var stage = Empty // Current stage of the exchange.

  /** A piece of data currently being exchanged. */
  private var slot = null.asInstanceOf[A]

  private val lock = new Lock

  /** Condition to signal to a thread waiting to exchange. */
  private val canStart = lock.newCondition

  /** Condition for the second thread to signal to the first thread that it
    * has swapped values. */
  private val swapped = lock.newCondition

  def exchange(x: A): A = lock.mutex{
    canStart.await(stage != Exchanging) // Wait for previous exchange to
                                     // finish (1).
    if(stage == Empty){
      slot = x; stage = Filled // Deposit my value. 
      canStart.signal() // Signal to second thread at (1).
      swapped.await() // Wait for second thread (2).
      assert(stage == Exchanging)
      stage = Empty; canStart.signal() // Signal to next first thread at (1).
      slot
    }
    else{
      assert(stage == Filled)
      val result = slot; slot = x // Swap first thread's value with mine.
      stage = Exchanging; swapped.signal() // Signal to first thread at (2).
      result
    }
  }
}
\end{scala}
\caption{An implementation of an exchanger using conditions.}
\label{fig:MonitorExchanger}
\end{figure}

%%%%% Check position of figure

\begin{figure}
\begin{center}
\def\height{5.7}
\begin{tikzpicture}[xscale = 2.8, yscale = 0.7]
\draw (0,0) node[draw] (t1) {Thread 1};
\draw (sender) -- ++ (0,-\height);
\draw (1,0) node[draw] (t2) {Thread 2}; 
\draw (receiver) -- ++ (0,-\height);
\draw (2,0) node[draw] (sender2) {Next thread 1};
\draw (sender2) -- ++ (0,-\height);
%
\draw[<-] (t1) ++ (0,-2) -- 
  node[above]{\scalashape\footnotesize canStart} ++ (-1,0); 
\draw[->] (t1) ++ (0,-3) -- 
  node[above]{\scalashape\footnotesize canStart} ++ (1,0);
\draw[->] (t2) ++ (0,-4) --
  node[above]{\scalashape\footnotesize swapped} ++ (-1,0);
\draw[->] (t1) ++ (0,-5) --
  node[above, near start]{\scalashape\footnotesize canStart} ++ (2,0);
\end{tikzpicture}
\end{center}
\caption{A sequence diagram illustrating the signals in the implementation of
  an exchanger.}
\label{fig:exchanger-seqDiagram}
\end{figure}

%%%%%

We use a |Lock| to provide mutual exclusion, together with two conditions
for signalling.  The condition |canStart| is used to prevent a third thread
interfering with an exchange: a signal on this condition indicates that that
thread can now start an exchange (unless $\sm{stage} = \sm{Exchanging}$,
i.e.~unless it has been preempted by another thread).  The condition |swapped|
is used by the second thread to signal that it has swapped values.
Figure~\ref{fig:exchanger-seqDiagram} illustates the sequence of signals: the
choreography of the exchange.

In the |exchange| operation, a thread starts by waiting, on |canStart|, for
the previous exchange to finish, i.e.~until $\sm{stage} \ne \sm{Exchanging}$. 

If $\sm{stage} = \sm{Empty}$, then this is the first thread.  It deposits its
value, and signals on |canStart|.  It then waits for a signal from the second
thread on |swapped|, indicating that it has swapped values (so $\sm{stage} =
\sm{Exchanging}$).  The thread resets for the next round, signals on
|canStart| to a waiting thread, and returns the value deposited by the second
thread.

Otherwise, in the case $\sm{stage} = \sm{Filled}$, this is the second thread.
It swaps its value with that deposited by the first thread, signals to the
first thread on |swapped|, and returns the first thread's value.

Note that the value of |stage| cycles between |Empty|, |Filled|, and
|Exchanging|.  Also, the signal on |swapped| guarantees that $\sm{stage} =
\sm{Exchanging}$: the assertion that captures this is useful documentation,
but would also help to capture some possible errors in an implementation. 

\begin{instruction}
Study the details of the implementation.
\end{instruction}
